\section{Patch extraction and cube allocation}
\label{sec:patch-tiling-under-9}

We enter this construction with the initial discrepancy
\eqref{eq:source-2}, the deep discrepancy \eqref{eq:source-9} in both
orientations, and the eventual cluster-witness absence of
Corollary~\ref{cor:eventual-cluster-absence}, all on the same retained
sides. The persistent discards in Corollary~\ref{cor:initial-discrepancy}
and the subsequent stabilizations of Lemma~\ref{lem:stabilization} total
\(o(N)\), so each retained side has size \(N-o(N)\). We will extract
disjoint host patches, select one common construction for them, allocate
cube roles by prefixes, and then clean their first supports for prescribed
comparison laws.

\subsection{Measured bias and cluster scales}
\label{sec:measured-patch-scales}

Use a tiny \(\iota>0\), say less than \(\min(x_*,\eta_0,.01)/1000\), and \(0<a_C<\min(\eta_0,1)/10^6\). Choose \(a_B>0\) much smaller than \(a_C\), also small enough that \eqref{eq:source-9} with exponent slack \(\iota/2\) applies at widths \(n^{a_B}\) on both sides. Reserve a disjoint part of each retained side of size, say, \(\lfloor N/3\rfloor\) for possible later use.

\begin{definition}[Residual scales]\label{def:residual-scales}
On the labels left after this reserve and any earlier patch removals, measure two dyadic scales (values that are powers of two with integer exponent). Both maxima are recalculated on each current residual pair:
\begin{itemize}
\item \(g\) is the largest dyadic value \(b\ge2\) for which two uniform sets each of size \(\ge N e^{-b^{a_B}}\) have absolute bias at least \(b/n\).
\item \(q\) is the largest dyadic \(b\ge2\) for which, in some orientation, there is a uniform first set of size \(\ge N e^{-b^{a_C}}\) and a family of disjoint bins on the second side, each of size at least \(e^b\), with union of size \(\ge N e^{-b^{a_C}}\), such that all distinct pairs in each bin have inner product of their \(f\)-vectors \(>\theta\) under that first law. This inner product does not depend on color choice.
\end{itemize}

Use 1 for an empty maximum.
\end{definition}

These scales measure actual uniform-set witnesses at the current dimension;
the exponents of Definition~\ref{def:limiting-bias-exponents} measure
availability along the stabilized sequence.

\begin{lemma}[Bounds for the residual scales]\label{lem:residual-scale-bounds}
Always \(g\le n^{\iota/2}\) eventually by its elementary cap \(n\) and the budget just chosen. Also \(q<n^\gamma\) eventually for \emph{each} fixed \(\gamma>0\), uniformly on the residual sets.
\end{lemma}

\begin{proof}
A nontrivial bias witness has \(g\le n\), so its two uniform laws have widths at most \(n^{a_B}\). The discrepancy estimate chosen with slack \(\iota/2\) gives
\[
 g/n\le n^{-1+\iota/2},
\]
which proves the first bound. For the cluster-scale bound, if not, pass to maxima with \(\log_n q\to\gamma'\in[\gamma,1]\) on a subsequence (\(q\le\log N\) when nontrivial). Remove column degree outliers into the stated first set using \eqref{eq:source-2}; this costs negligible mass even relative to the union size (exception count at most \(2Ne^{-n^{\eta_0}}\), testing their uniform laws). Here we can take deviation \(2n^{-\eta_0}=o(1)\), since the first width exponent is at most \(\gamma' a_C+o(1)\). Keep the bins retaining at least half their mass, restricted to the survivors.

For distinct retained labels, their codegree in the red color is
\[
 \frac{1+m(y)+m(y')+K(y,y')}{4},
\]
where the means and inner product are under the first law. Here \(|m(y)|,|m(y')|\le4n^{-\eta_0}\) and \(K(y,y')>\theta\); blue changes the signs of the two means and keeps \(K\). Thus both codegrees exceed \(1/4+\theta/8\) eventually. A diagonal codegree is the degree itself, close to \(1/2\). All pairs therefore have codegree \(\ge1/4+n^{-\delta}\), where one can fix rational numbers \(\zeta,\delta\) with \(0<\zeta<\gamma'\) and \(\gamma'a_C<\delta<\min(\eta_0,\zeta,1)/2000\). Uniform bin clusters with weights proportional to retained sizes now give a full-dimensional cluster witness with atom cap \(\exp(-n^\zeta)\) and both widths \(\le n^\delta\), contradicting Corollary~\ref{cor:eventual-cluster-absence}.
\end{proof}

\subsection{Extraction and cube allocation}
\label{sec:extraction-and-cube-allocation}

Choose the parameters in the following order. There is room in the
inequalities, and very large \emph{fixed} thresholds suffice.
\begin{enumerate}
\item For extraction, use a large constant \(M_1\) compared to the conflict-count constant \(C_*\) above, so \(2C_*/\sqrt{M_1}\ll 10^{-3}\). Let \(C_b>100 a_C/a_B+100\), \(M_{\rm lo}>C_b+100\), \(0<c_q<1/(20M_{\rm lo})\), and \(M_{\rm hi}\gg M_{\rm lo}+10/c_q\). Fix tiny \(\omega>0\) with \(5\omega M_{\rm hi}<a_C/100\). Here \(M_1\) ensures that the homogeneous extension count in direct mode fits its gain; \(C_b\) separates the larger bias-test width from the input-law width. The constants \(M_{\rm lo},M_{\rm hi}\) set internal dimensions, and \(\omega\) keeps tuple-filter costs below the available width and bin-size budgets.
\item For cluster patches, use the fixed gain parameter \(a=\theta/100\). For the geometric constructions take a tiny \(\rho>0\) so that \(h\)-bit balls of radii up to \(1000\rho h\) have log volume \(< (a/10^9)h\) for large \(h\).
\item Use large constants \(K_B,A_0\) compared to \(10^6R\); \(A_0\) serves for order-\(\log n\) late classes. Fix all the small-surplus geometric and probability constants before taking a large scale threshold \(Q_0\). Increase \(Q_0\) as needed so estimates in positive powers of \(q,g,h\) below already have the required slack for arguments bounded below by \(Q_0\). They need not all grow along the sequence.
\end{enumerate}

For a patch \((X_i,Y_i)\), write \(M_i=|X_i|=|Y_i|\). Its
internal dimension \(h_i\) is the number of cube directions to be handled
inside an internal slice; \(s_i\) is the gain budget for support and
prefix losses. The following three patch types specify the output of one
extraction. Their scales \(g,q\) are those of the residual pair at that
pass, and the constants in the size bounds are independent of \(n\).

\begin{itemize}
\item In the \emph{bounded case}, \(\max(g,q)<M_1Q_0\). We use one
patch of size \(M_i\asymp N\), with uniform laws on its two sides,
either orientation and color, and \(h_i=s_i=0\). It belongs to low
direct mode in the later assignment.
\item A \emph{direct bias patch} is used when the bounded case does not
apply and \(g>M_1q\). It has
\[
 \begin{gathered}
 M_i\ge cNe^{-g^{a_B}},\qquad
 d_G(x,U(Y_i))\ge\tfrac12+\frac{g}{4n}\quad(x\in X_i),\\
 h_i=0,\qquad s_i=10^{-3}g,
 \end{gathered}
\]
where \(U(Y_i)\) is uniform on \(Y_i\). It is high when
\(g>K_B\log n\), and low otherwise.
\item A \emph{cluster patch} is used in the remaining case
\(g\le M_1q\), where \(q\ge Q_0\). It has
\(M_i\ge cNe^{-q^{a_C}}\), and \(Y_i\) is partitioned into bins
of common size \(d_i\). Every pair in a bin, including repeated labels,
has \(G\)-codegree at least \(1/4+3a\) under the uniform law on \(X_i\).
It is low when \(q\le(\log n)^{c_q}\), and high otherwise. High patches
use high-small-bin mode when \(q\le(\log n)^2\), and high-large-bin mode
otherwise. Set
\[
 d_i=\lfloor e^{q/2}\rfloor,
\]
except in high-small-bin mode, where this is replaced by the smaller of
\(\lfloor e^{q/2}\rfloor\) and
\(\lfloor\exp(\sqrt{\log n})\rfloor\). Choose a power of two
\(q^M\le h_i<2q^M\), taking \(M=M_{\rm lo}\) in low mode and
\(M=M_{\rm hi}\) in high mode, and put
\[
 k_i=\lceil h_i^{3\omega}\rceil,\qquad
 T_i=\lceil h_i^\omega\rceil,\qquad s_i=(a/10^6)h_i.
\]
The internal solver of Section~\ref{sec:internal-solver-and-profiles-for-these-tilings}
will turn the within-bin codegree surplus into a gain on the \(h_i\)
internal directions.
\end{itemize}

In the proposition, \(S\) denotes the total size of the selected
mode/orientation/color family before a subcollection is retained; the
dyadic sum and prefix claims concern that allocated subcollection.

\begin{proposition}[Extraction and allocation]\label{prop:patch-extraction-allocation}
There are pairwise disjoint patches \((X_i,Y_i)\) with \(|X_i|=|Y_i|=M_i\) and \(\sum_i M_i=S\ge c\,N\) for a fixed \(c>0\), all of one mode, orientation and color. They have one of the three patch types just specified. A subcollection admits dyadic prefix lengths \(\ell_i\) with
\[
 \sum_i2^{-\ell_i}=1,\qquad
 M_i/S\le 2^{-\ell_i}<2M_i/S.
\]
In its prefix allocation, patch \(i\) has \(\Theta(2^nM_i/N)\) roles of each parity. Each star in that patch has exactly \(\ell_i\) crossing flips, visiting distinct other leaves, and a nested tail coordinate set of size \(h_i\) supplies its internal directions. Eventually \(\max_i(h_i,\ell_i)<n^\iota\). Except in the bounded case,
\[
 \ell_i,\ \log(N/M_i)\le s_i/(1000u).
\]
The bounded case uses one patch, with \(\ell_i=h_i=s_i=0\).
\end{proposition}

\begin{proof}\leavevmode\par
At every pass, work on the labels left after removing both supports of the
earlier patches from the working sides, and recompute \(g,q\) on the current
residual pair. A new pass is
started only while the extracted equal sizes sum to less than \(.1N\).
Thus each residual side then has size at least
\[
N-o(N)-\lfloor N/3\rfloor-.1N\ge cN
\]
for a fixed \(c>0\), as needed by the bounded branch.

\paragraph{Extracting one patch.}
If \(\max(g,q)<M_1 Q_0\) during extraction, take the bounded case right away on those residual sets, still of sizes \(\ge cN\); this uses just one equal-side-size patch, by truncating to arbitrary sets of the smaller size, with uniform laws, either orientation and color, and low direct mode below. Otherwise proceed as follows (later \(g,q\) for an extracted patch refer to the values here at its own extraction):
\begin{itemize}
\item If \(g>M_1 q\), use direct bias mode. Extract first and second sets \(X_i,Y_i\) of equal size \(M_i\ge c N e^{-g^{a_B}}\), with the second law uniform and every first label having degree at least \(1/2+g/(4n)\) into it in some fixed color of that patch. To see this, let \(U,V\) witness \(g\) and choose their denser color. The top eighth \(U_1\) of rows by degree has size at least \(Ne^{-(2g)^{a_B}}\), since the fixed threshold makes \((2g)^{a_B}-g^{a_B}>\log8+O(1)\). The absent \(2g\) budget therefore gives \(d_G(U_1,V)<1/2+2g/n\). If fewer than an eighth of all rows had degree at least \(1/2+g/(2n)\), the remaining seven eighths would be below that threshold, and
\[
 d_G(U,V)<\tfrac18(1/2+2g/n)+\tfrac78(1/2+g/(2n))
             =1/2+11g/(16n),
\]
contrary to the witness. Let \(U^+\) be the rows meeting this degree threshold and take \(M_i=\min(|U^+|,|V|)\). Keep any \(M_i\) of the good first labels, and choose an \(M_i\)-subset of \(V\). A uniform such choice preserves all their degrees to \(o(1/n)\) simultaneously: \(M_i=e^{\Omega(n)}\), so sampling concentration permits a union over at most \(N\) rows. One realization gives the asserted \(g/(4n)\) surplus.
\item Otherwise use the cluster mode (\(q\ge Q_0\)), with orientation from its witness and, say, red color. Extract \(X_i,Y_i\) of equal size \(M_i\ge cN e^{-q^{a_C}}\), where \(Y_i\) is partitioned into bins of the prescribed size \(d_i\). Every within-bin pair (including diagonals) has codegree \(\ge1/4+3a\) under uniform \(X_i\). First remove the degree outliers as in Lemma~\ref{lem:residual-scale-bounds}. Retain original bins losing at most half their mass, restrict them to the survivors, and partition them into bins of the stated size, discarding the final remainders. This loses only a negligible or fixed small fraction. Choose a common size that is a multiple of \(d_i\), no larger than either the surviving first set or the whole-bin union, and maximal with these properties. It remains at least \(cNe^{-q^{a_C}}\). Keep that many whole-bin labels on the second side and sample the same number of first labels. Since this number is \(e^{\Omega(n)}\), concentration followed by a union over the at most \(N^2\) within-bin pairs preserves every codegree with the displayed slack.
\end{itemize}

The chosen cluster parameters satisfy
\[
 k_iT_i=O(q^{4\omega M})\ll q^{a_C/2},\ \log d_i.
\]
For high-small-bin mode, \(\log d_i=\min(q/2,\sqrt{\log n})+O(1)\); the bound \(q\le(\log n)^2\) and the choice of \(\omega\) give the second comparison also at the cutoff. Likewise \(h_i\) is smaller than any fixed positive power of \(d_i\) needed below, uniformly after increasing \(Q_0\). The inequalities \(c_qM_{\rm lo}<1/20\) and \(c_qM_{\rm hi}>5\) give, respectively, \(h_i\le(\log n)^{.1}\) in low mode and \(h_i>(\log n)^5\) in high mode eventually. After either nonbounded extraction, remove its two supports before the next pass and retain its measured \(g,q\) values with that patch.

\paragraph{Repeating extraction and selecting one type.}
Continue until the sum of the equal sizes first reaches \(.1N\), or until
the immediate bounded branch fires. In the former case, select one of the
finite mode/orientation/color possibilities, distinguishing also
high-small-bin and high-large-bin. Let \(I\) be its patch indices and put
\(\sum_{i\in I}M_i=S\ge c'N\). In the immediate bounded case, ignore
earlier extractions and take \(I\) to contain only the bounded patch, with
\(S=M_i\) and \(s_i=h_i=0\). From now on all patches used are of this one
type, orienting notation accordingly. The color is selected together with
the mode and orientation, before any cube allocation; thus every edge
required by the subsequent construction is tested in this single selected
color.

\paragraph{Dyadic prefix allocation.}
For \(i\in I\), put \(p_i=M_i/S\) and round it \textbf{up} to a dyadic
weight \(r_i=2^{-\ell_i}\). Then
\[
p_i\le r_i<2p_i,\qquad \sum_{i\in I}r_i\ge1.
\]
Order the \(r_i\) decreasingly, and let \(J\) be the first segment whose
sum reaches \(1\). Both \(1\) and all preceding terms are integer multiples
of the last term in this segment. The preceding partial sum is below \(1\),
so adding that last term reaches \(1\) exactly. Hence
\(\sum_{i\in J}r_i=1\). We retain only these patches for the allocation,
but \(S\) continues to mean the total over \(I\) before this restriction.
In particular \(\sum_{i\in J}M_i>S/2\).

The allocation uses subcubes obtained by fixing initial coordinates, as in
the cube tilings of Conlon, Fox, Lee and
Sudakov~\cite[Section~2.2]{conlon-fox-lee-sudakov2016} and their adaptation
by Fiz Pontiveros, Griffiths, Morris, Saxton and
Skokan~\cite[Section~2]{fiz-pontiveros-griffiths-morris-saxton-skokan2016}.
Here the sizes are chosen from the patch masses, and the later embedding
uses the patch laws. Use the weights \(r_i\), \(i\in J\), as disjoint
complete prefix leaves in the first cube coordinates. Positions in leaf
\(i\) belong to patch \(i\) for both parities. The leaf has exactly
\(2^{n-\ell_i-1}=\Theta(2^nM_i/N)\) roles of each parity eventually.
Leading flips, called crossings, number \(\ell_i\) per star of patch
\(i\), visiting distinct other leaves: the first prefix mismatch identifies
the flipped coordinate.

\paragraph{Internal and bulk coordinates.}
In the tail, beyond all leading coordinates, dedicate nested coordinate
sets of sizes \(h_i\) as internal for patch \(i\); by a slice in that patch
we mean exactly fixing their complementary coordinates. Other nonleading,
noninternal flips there are called bulk. Here nonleading is relative to the
current leaf: a coordinate beyond its prefix may still belong to a longer
leaf's prefix. The size bounds and dyadic rounding give
\[
 \log(N/M_i)\le
 \begin{cases}g^{a_B}+O(1)&\text{direct},\\q^{a_C}+O(1)&\text{cluster},\end{cases}
 \qquad \ell_i\le\log_2(S/M_i)+1.
\]
Together with \(g\le n^{\iota/2}\), and Lemma~\ref{lem:residual-scale-bounds} used with a fixed \(\gamma<\iota/(2M_{\rm hi})\), these show that the maximum of the \(h_i\) over \(I\) and the \(\ell_i\) over \(J\) is \(<n^\iota\) eventually. In particular all asserted
dimension allocations fit. Except in the bounded case we can ensure
\[
 \ell_i,\ \log(N/M_i)\ \le\ s_i/(1000u)
\]
because \(g^{a_B}=o(g)\) and \(q^{a_C}=o(q^M)\) as their arguments grow; choose the fixed threshold so these inequalities hold with the displayed constants for every argument above it. In the bounded case these quantities are \(O(1)\) (actually \(\ell_i=0\)).
\end{proof}

\subsection{Cleaning at a prescribed input profile}
\label{sec:cleaning-at-input-profile}

Section~\ref{sec:internal-solver-and-profiles-for-these-tilings} will choose
the cluster comparison laws through a fixed point. We first clean the
first supports at any prescribed input, with enough stability for nearby
inputs. Here \(X_i\) continues to denote the extracted set of size \(M_i\);
\(\mu_i\) is uniform on the support left after cleaning.

\begin{lemma}[Stable cleaning]\label{lem:stable-patch-cleaning}
Here are the cleaned first laws. In direct and bounded modes the comparison laws \(\pi_i\) are uniform on \(Y_i\). In cluster mode for now consider any vector of input comparison laws \(\pi_i\) on \(Y_i\) of normalized cap
\[
 M_i\max\pi_i\le \exp(10 k_iT_i).
\]
At such input, take \(\mu_i\) uniform on \(X_i\) after the following cleanings, which need waste at most a tiny fraction (less than \(a\)). Its support has
\begin{equation}
 \begin{array}{ll}
  1/2+g/(4n)\le D_{\pi_i}\le1/2+4g/n &\text{direct bias},\\
  |D_{\pi_i}-.5|\le 10q^{C_b}/n &\text{cluster mode},\\
  |D_{\pi_i}-.5|\le K/n &\text{bounded case, fixed }K.
 \end{array} \label{eq:source-12}
\end{equation}
In cluster mode the right deviation budget \(10q^{C_b}\) can be made smaller than \(s_i/(100u)\). Furthermore, for a parameter \(Q_i\), each cleaned support has no set of \(\lceil e^{Q_i}\rceil\) distinct labels with \(K_{\pi_i}(x,z)>3\theta\) on every pair, as required by Lemma~\ref{lem:homogeneous-peeling}. This parameter satisfies \(\Delta_i:=\exp(C_* Q_i)\le e^{s_i}\), or just \(\Delta_i\le K\) fixed in the bounded case. All first labels there have \(|D_{\pi_j}-.5|\le 3b_*\) for \(j\ne i\), and \eqref{eq:source-11} with its relaxed thresholds \(n^{-1-.02},2\xi,100n,e^{-n^{.19}}\), with \(z\) uniform on the original extracted \(X_i\), for every patch law \(\pi_j\). Cluster mode retains all within-bin codegrees \(\ge1/4+a\) into \(\mu_i\). In cluster mode the same support has these properties for all inputs in a sufficiently small \(\ell^1\)-neighborhood of the input used for cleaning. The no-clique assertion is about each such support separately, not their union over inputs.
\end{lemma}

\begin{proof}\leavevmode\par
\paragraph{Approximating a capped law by a uniform set.}
The measured budgets concern uniform sets, so first fix a law \(\pi\) on the residual labels of width \(W=\log(N\max\pi)\), and a larger powered width budget \(B\). Set
\[
 t=Ne^{-(B+W)/2},\qquad \Pr(y\in V)=t\pi(y),
\]
independently over labels. Then \(t\max\pi=e^{-(B-W)/2}<1\) and \(\mathbb E|V|=t\gg Ne^{-B}\). In the applications \(B,W=o(n)\), so \(t=e^{\Omega(n)}\). Concentration of the cardinality and all bounded degree and pair sums, followed by a union over \(O(N^2)\) tests, gives one realization whose uniform law approximates every such test against \(\pi\) to \(o(1/n)\), and whose size is at least \(Ne^{-B}\).

\paragraph{Own-patch degrees.}
At the residual pair used to extract a cluster patch, choose a dyadic bias budget in \([q^{C_b},4q^{C_b}]\), larger than the measured \(g\le M_1q\). Its powered width exceeds the input width with ample slack:
\[
 q^{C_ba_B}\gg q^{a_C}+q^{4\omega M},\qquad
 \log(N\max\pi_i)\le\log(N/M_i)+10k_iT_i.
\]
Use the uniform approximant just constructed for the second law. If the rows of \(X_i\) exceeding \(8q^{C_b}/n\) in either bias direction numbered at least the higher budget's size threshold, they and this approximant would form a rectangle witnessing that absent budget. Each exceptional row set is therefore smaller than its threshold, a tiny fraction of \(M_i\). In direct mode \(\pi_i\) is already uniform and fits the absent \(2g\) budget; the same rectangle test gives the upper degree trimming while preserving the lower bound. In bounded mode choose fixed dyadic budgets above \(M_1Q_0\), large enough that their size thresholds are less than \(aM_i/10\). They give a fixed degree constant \(K\).

\paragraph{Own-patch clique exclusion.}
For cluster mode choose a dyadic \(Q_i\in[q^2,2q^2]\); for direct mode choose one in \((g/(2\sqrt{M_1}),2g/\sqrt{M_1}]\). Both exceed the measured cluster maximum. Their powered widths exceed the input widths because
\[
 q^{2a_C}\gg q^{a_C}+k_iT_i,
 \qquad (g/\sqrt{M_1})^{a_C}\gg g^{a_B}
\]
above the fixed threshold. Approximate \(\pi_i\) by a uniform law at this budget. Greedily remove disjoint cliques of size \(\lceil e^{Q_i}\rceil\) in \(X_i\), with pair inner products under \(\pi_i\) greater than \(2\theta\). If their union had size at least \(Ne^{-Q_i^{a_C}}\), it would, together with the approximating first law in the reverse orientation, witness the absent cluster budget \(Q_i\): the approximation retains pair inner products greater than \(\theta\). The removed union is thus smaller than that threshold. Its fraction is tiny, and \(\Delta_i=e^{C_*Q_i}\le e^{s_i}\) by the powers of \(q\) and the choice of \(M_1,Q_0\). In bounded mode a sufficiently large fixed \(Q_i\) gives the same small-loss conclusion and a fixed \(\Delta_i\).

\paragraph{Other-patch degrees and row tails.}
In cluster mode the input widths satisfy
\[
 \log(N/M_i)+10k_iT_i\le n^{x_*/4},
 \qquad \#I\le e^{n^{2\iota}}=e^{o(n^{.2})}.
\]
For direct and bounded modes omit the term \(10k_iT_i\); their input laws are uniform. For each \(j\ne i\), conditioning uniform \(X_i\) on a putative outlier fraction \(e^{-cn}\), with fixed \(c<\alpha/2\), still fits the linear width in \eqref{eq:source-9}. Thus the fraction with deviation greater than \(2b_*\) against \(\pi_j\) is \(e^{-\Omega(n)}\). Lemma~\ref{lem:simultaneous-row-trimming} supplies the row-tail removal for the whole collection. Their unions have vanishing relative mass. Choose the fixed thresholds so own-degree and clique removals together cost less than, say, \(a/2\); the other two removals together cost less than \(a/2\) eventually. If a total fraction \(r<a\) is removed, each original within-bin codegree of at least \(1/4+3a\) becomes at least
\[
 \frac{1/4+3a-r}{1-r}\ge1/4+a.
\]

Finally hold this cleaned support fixed while varying the inputs. Degrees and pair inner products are Lipschitz in the \(\ell^1\) distance of the comparison laws. A sufficiently fine mesh, of diameter at most \(n^{-3}\), preserves the degree margins \(8\) to \(10\), \(2b_*\) to \(3b_*\), the clique threshold \(2\theta\) to \(3\theta\), and the relaxed row-tail bounds proved in Lemma~\ref{lem:simultaneous-row-trimming}. Cleaning is performed separately at each corner; it requires no union over mesh corners. The later union of corner supports inherits only the bounds that hold label by label.
\end{proof}

The threshold \(Q_0\) is fixed before the asymptotic sequence is taken. In particular, estimates in powers of a cluster scale are required uniformly above this threshold; low-mode scales are not assumed to diverge. For profile interpolation, clique exclusion is a property of each cleaned corner support separately. The later union of those supports is used only for the degree and row-tail bounds that hold label by label.
